\section{总结与展望}

本文围绕“多视图几何重建中的学习方法”这一主题，以 VGG 团队为例梳理了近五年的代表性研究主线。综合全文可以看到，VGG 团队的重要性并不只在于若干单篇论文分别提升了某项局部性能，而在于其成果之间形成了一条较完整、可解释的研究链：先以几何先验、数据基础与任务边界回答“学习方法能否真正进入几何问题”，再以点轨迹和对应学习回答“系统应当依赖什么样的观测”，随后以可微 SfM 回答“几何求解如何被重组为可学习系统”，最后以 VGGT 等工作进一步追问“统一表示能否吸收更多后端逻辑”\cite{wu2020unsup3d,wang2021deep2view,reizenstein2021co3d,bhalgat2023lighttouch,karaev2024cotracker,wang2024vggsfm,wang2025vggt}。因此，VGG 团队的研究不是若干论文的并列堆叠，而是一条由问题设定、观测组织、系统求解到表示统一逐步展开的清晰主线。

从研究特色上看，VGG 团队至少表现出四方面鲜明特点。第一，它始终以多视图几何中的可解结构为中心，而不是把三维重建理解为一般性的黑盒回归任务。第二，它强调系统级推进，不满足于只优化某个局部模块，而是不断重新组织观测、求解与表示之间的关系。第三，它重视真实数据、伪标签与基准构建，使研究主线不仅有方法创新，也有训练和评测基础设施的持续支撑\cite{reizenstein2021co3d,karaev2025cotracker3}。第四，它对统一几何表示保持持续关注，使模型输出不只服务于单次重建，而是逐步具备跨任务复用的潜力\cite{wang2025vggt,sucar2025dpm}。这些特点共同决定了 VGG 团队在该领域中的位置：它既不是完全站在传统 SfM 一侧，也不是简单站在“去几何化”的纯前馈路线一侧，而是处在“保留几何结构、同时推动学习化系统重构”的关键交汇点。

从对领域发展的推动作用看，VGG 团队的贡献主要体现在三个层面。其一，在观测层面，它推动研究者重新理解高质量对应的含义，使对应关系从脆弱的成对匹配逐步转向更稳定、更密集、可跨时空延续的轨迹表示\cite{bhalgat2023lighttouch,karaev2023dynamicstereo,karaev2024cotracker}。其二，在系统层面，它通过 VGGSfM 证明传统结构恢复并非只能依赖分散的不可微模块，而可以被重写为端到端可训练、可整体优化的几何系统\cite{wang2024vggsfm}。其三，在表示层面，它通过 VGGT 与 Dynamic Point Maps 等工作推动多视图几何从“优化驱动的流水线”进一步走向“表示驱动的统一建模框架”\cite{wang2025vggt,sucar2025dpm}。换言之，VGG 团队推动的不只是某个 benchmark 指标的上升，更是整个方向从模块增强走向系统重构、再走向统一表示的范式迁移。也正是在这个意义上，VGG 团队可以被视为当前多视图几何学习化进程中的代表性团队之一。

如果给出本文对其研究体系的整体评价，我更倾向于把 VGG 团队理解为当前多视图几何学习化进程中的“代表性中轴”。其优势在于推进路径稳健、问题意识清楚、前后工作衔接自然，因而特别适合用来观察这一方向如何从模块增强逐渐走向系统重构与统一表示。更重要的是，这一研究体系并没有脱离几何问题本身，而是在每一步都围绕“几何结构如何被保留、吸收和重用”这一核心矛盾展开。与此同时，这条主线也保留了值得继续警惕的张力：统一前馈模型虽然提高了效率与接口一致性，但也显著增强了对训练数据规模、监督质量和模型容量的依赖；而几何可解释性虽然仍被保留为重要目标，却尚未在动态、长尾和复杂开放场景中彻底转化为稳定优势\cite{wang2024vggsfm,wang2025vggt}。因此，对 VGG 主线最准确的评价不是“已经给出最终答案”，而是“较清楚地指明了下一阶段问题会出现在哪里”。

结合全文分析，本文最终想强调的综合结论是：多视图几何重建中的学习方法并没有削弱几何本身的重要性，反而促使几何约束以更灵活、更深层的形式进入特征学习、监督构造、系统设计与统一表示之中。VGG 团队之所以值得被单独综述，也正是因为它较完整地展示了这一变化过程。面向后续发展，可以预见该方向仍会围绕三个问题继续推进：其一，如何在减少显式优化的同时保持高精度和强可解释性；其二，如何利用更大规模真实数据与弱监督信号提升统一模型的泛化能力；其三，如何把当前以静态场景为主的统一几何表示进一步扩展到动态场景、开放世界感知与三维内容生成等更复杂任务中。总体而言，VGG 团队的研究体系尚未封闭，但其已经为这一方向的下一阶段演进给出了较清晰的路径指示。
